\section{Introduction}

Countries that consume more sugar are fatter. Plot national sugar supply against adult obesity prevalence for 37 Sub-Saharan African countries and the correlation is strong: $r = 0.50$ after controlling for GDP, urbanization, and vegetable oil supply ($p = 0.002$). Drop any single country and the result holds (leave-one-out range: 0.42--0.63). Add subregion fixed effects and the correlation strengthens to $r = 0.63$. By every conventional cross-sectional standard, sugar supply predicts obesity.

The association is a between-country development gradient. When we follow the same 37 countries over 13 years, tracking whether rising or falling sugar supply corresponds to faster or slower obesity growth, the sugar coefficient collapses toward zero. Country fixed effects reduce it to $t = 0.25$. Adding year effects gives $t = -0.64$. Detrending each country's individual trajectory gives $t = -0.13$. Three further approaches reach the same boundary: first differences ($r = -0.03$), long differences over the full decade ($r = -0.19$), and a Mundlak correlated random effects model ($t = -0.66$). Six specifications, each removing a different class of confounding, locate the signal between countries rather than within them.

The explanation is structural. Obesity has risen along near-linear trajectories in virtually every country over the past decade. Entity and year fixed effects together absorb 99.4\% of obesity variance in our panel. Food supply variables compete for the remaining 0.6\%, a residual space dominated by measurement error and short-run noise. The same boundary appears across all 10 FAOSTAT food groups and across a targeted grand-total energy, fat, and protein comparator under two-way fixed effects and first differences.

The mechanism is visible in the estimand. Cross-sectional studies compare countries at a point in time, so they combine persistent between-country differences with the shared upward trend in obesity. A correlated random effects decomposition confirms this: the between-country sugar coefficient is $t = 3.46$ while the within-country coefficient is $t = -0.66$. The association lives between countries. A soybean oil case study makes the mechanism concrete. Soybean oil supply passes the same cross-sectional and one-way fixed-effects tests that sugar passes, yet its lead-lag ratio is 0.99: pure co-trending with no temporal direction.

The paradox matters because cross-sectional food supply studies have shaped policy. Basu et al.\ \cite{basu2013} applied generalized estimating equations to repeated cross-sections in 175 countries and argued that sugar availability independently predicted diabetes prevalence. A companion study \cite{basu2013ajph} linked soft drink consumption to overweight across 75 countries using purely cross-sectional regression, a finding cited over 500 times and invoked in WHO guidance on sugar-sweetened beverage taxation \cite{who2023ssb}. Individual-level evidence on sugar-sweetened beverages remains compatible with our result; the failure lies in national food-supply identification. Both Basu studies draw their identification from between-country variation, exactly the variation we show is confounded by shared trends. More broadly, the Global Burden of Disease project assigns diet-attributable deaths using food supply data that faces the same trend-domination problem \cite{gbd2019}, and a recent global analysis attributed 2.2 million new diabetes cases to sugar-sweetened beverages using dietary intake estimates informed in part by national food supply data \cite{laracastor2025}. Cross-national food-supply designs inherit the confounding structure we document when they use between-country variation to identify dietary effects.

The same trend structure appears worldwide. Among 200 countries in the NCD-RisC database \cite{ncdrisc2024bmi}, 98\% show linear obesity trends with $R^2 \geq 0.90$ over 2010--2021. This near-universal linearity means that any food supply variable trending over the same period will correlate with obesity cross-sectionally, regardless of whether it causes obesity.

Initial obesity level predicts subsequent obesity change ($t = 3.45$, $R^2 = 0.41$). Countries that started with higher obesity in 2010 gained more over the following decade. This divergence pattern points to structural economic and demographic conditions as the driver of cross-country obesity differences.
